\documentclass[12pt]{article}

\usepackage{graphicx}
\usepackage{amsmath,amsthm,amsfonts}                % AMS Math
\usepackage{thmtools}                                       % Theorem Tools
\usepackage{bm}                                             % Bold Math

\title{LaTex test kram}
\author{Testautor}
\date{April 2022}

\begin{document}

\maketitle

% Start vom Fließtext
\pagebreak
\section{Einleitung}
Das ist ein test Text. Hier steht eine Formel direkt im Paragraphen: 
$\sum_{n = 1}^{\infty} x^n$. Hier steht dann wiederum eine Formel einzeln als Figure:
\begin{align}
  \int_{0}^{\infty} 6x^4-\frac{4}{6} \,dx \label{random integral}
\end{align}
Eine Formel kann man anhand ihres Labels referenzieren, dann passt sich
die Zahl automatisch an. Z.B. das integral (\ref{random integral}) wird automatisch
erkannt. \linebreak

% Tony Bild
\begin{figure}[h]
  \centering
  \includegraphics[width=0.25\textwidth]{Bilder/Unvernichtet.jpg}
  \caption{Ein schönes Bild}
  \label{fig:mesh1}
\end{figure}

% Liste
\begin{itemize}
  \item Das ist ein Listenpunkt!
  \item Das ist ein weiterer Listenpunkt!!
\end{itemize}
% Ein Enter ist egal, wird gedruckt als wäre es in einer Zeile
Das steht alles in
einer Zeile

%Eine leere Zeile startet einen neuen Paragraphen
Das ist ein neuer Paragraph.

Ich möchte jetzt hier einmal einen Testtext anfügen, damit wir mal sehen ob das mit der synchronisation klappt.
\pagebreak

Die Tabelle \ref{table:data} ist eine vollformatierte Tabelle. Toll!
\begin{table}[h!]
  \centering
  \begin{tabular}{||c c c c||} 
  \hline
  Col1 & Col2 & Col2 & Col3 \\ [0.5ex] 
  \hline\hline
  1 & 6 & 87837 & 787 \\ 
  2 & 7 & 78 & 5415 \\
  3 & 545 & 778 & 7507 \\
  4 & 545 & 18744 & 7560 \\
  5 & 88 & 788 & 6344 \\ [1ex] 
  \hline
  \end{tabular}
  \caption{Eine tolle Tabelle}
  \label{table:data}
\end{table}
\end{document}